\begin{claim}{D.2.27}{Proved}{The unsuspended first-cone functional vanishes with one common normalization}
Let $B=K(\FF_2,2)$, $A=Q\Sigma B$, and $Y=\Sigma^2B$.
Write $\alpha\in H^4(Y;\FF_2)$ for the suspended fundamental
class and $\beta=\operatorname{Sq}^1\alpha$.
There is one choice of the common primitive $c_8$ in
\textup{D.2.26}, with $dc_8=a_4\smile b_5$, for which
\[
 \theta_h=0\in H^8(A;\FF_2)
 \qquad\text{for every pointed map }h:A\longrightarrow Y.
 \tag{D.2.27a}
\]
Consequently, after any small simultaneous family of such
unsuspended attachments, a lift $\widehat\alpha$ to the common
cone satisfies
$\widehat\alpha\operatorname{Sq}^1\widehat\alpha=0$.
In particular the old degree-seven loop commutator survives.
This statement does not address exponents one or four, or
propagation through subsequent cones.

More generally, the same conclusion holds for one small
simultaneous family of maps $\Sigma^{k_i}A\to Y$ with
\[
 k_i\in\{0,2,3\}\cup\{5,6,7,\ldots\}.
 \tag{D.2.27a'}
\]
All these exponents use the same common primitive $c_8$;
the assertion permits mixtures of the indicated exponents.

The proof supplies an actual delooped relation. Put
$Z=\operatorname{fib}(Y\to K(\FF_2,4))$, using $\alpha$,
and $F=\Omega^2Z$. There are unique classes
$\vartheta_8\in H^8Z$, $\lambda_{10}\in H^{10}Z$ whose
twofold cohomology suspensions are $e_6^*,J_1^*$ of
\textup{D.2.24}. For the bottom $v_6\in H^6Z$,
\[
 \operatorname{Sq}^4\vartheta_8+
 \operatorname{Sq}^2\lambda_{10}
       =\epsilon v_6^2\quad(\epsilon\in\FF_2).
 \tag{D.2.27b}
\]
The common choice in \textup{(D.2.27a)} makes the universal
relative cup class on $Z$ exactly $\vartheta_8$.
\end{claim}
\begin{proof}
All cohomology and homology in this proof has coefficients
$\FF_2$. We first justify every delooping interface; no
cohomology-suspension isomorphism outside its indicated range
is assumed.

By the preceding low homotopy calculations,
$\pi_4Y=\ZZ/2$ and $\pi_5Y=0$. Thus $Z$ is five-connected,
$W=\Omega Z$ is four-connected, and $F=\Omega W$ is the
three-connected fiber already used in \textup{D.2.18--D.2.25}.
Moreover $H^6Z$ is one-dimensional because
$\pi_6Z=\pi_6Y=\ZZ/2$.

Consider the path fibrations for $Z$ and $W$.
Their Serre spectral sequences give the following low-degree
statements directly. Cohomology suspension identifies $H^8Z$
with $H^6F$. It also identifies $H^{10}Z$ with $H^9W$.
For $H^9W\to H^8F$, the only differential before transgression
out of the fiber group is
\[
 d_5:H^8F\longrightarrow H^5W\otimes H^4F.
 \tag{D.2.27c}
\]
The loop action on the path fibration identifies its dual
with multiplication of the transgressed bottom class by the
bottom fiber class. Explicitly the bottom tensor maps to
the Pontryagin product $A_4^2=W_8\in H_8F$.
This is the usual module property of the path-fibration
differential: the differential of a transgressive base
class tensored with a fiber cycle is its transgression
multiplied by that cycle.
Thus the kernel of \textup{(D.2.27c)} is exactly the
$W_8$-annihilator. The targets of $d_6,d_7,d_8$ vanish by
connectivity, and $d_9$ identifies this kernel with $H^9W$;
there is no earlier mixed source entering that base group.
Since $J_1^*$ vanishes on $W_8$, it has a unique lift to
$H^9W$ and then to $H^{10}Z$. This gives $\lambda_{10}$.
The lower isomorphism gives $\vartheta_8$.

For clarity, the kernel needed in degree twelve is equally
controlled. Before transgression into $H^{12}Z$, the only
possible mixed source in its path spectral sequence has
bidegree $(6,5)$; its image is $H^6Z\cdot H^6Z$.
Before transgression into $H^{11}W$, the only mixed sources
have bidegrees $(6,4),(5,5)$; their images lie in
$H^5W\cdot H^6W$. Cohomology suspension kills products,
so these are the exact kernels of the two respective
suspensions in these degrees. The image of
$H^{12}Z\to H^{11}W$ is primitive for loop multiplication.
Every degree-five and degree-six class of $W$ is primitive
by connectivity, and the reduced coproduct on their mixed
products is
\[
 x_5y_6\longmapsto x_5\otimes y_6+y_6\otimes x_5.
 \tag{D.2.27d}
\]
It is injective on the mixed-product subspace: the two tensor
bidegrees differ, so a zero coproduct forces the corresponding
tensor in $H^5W\otimes H^6W$ to vanish and hence its product
to vanish. Thus the second kernel has zero intersection with
the first suspension image. We have proved
\[
 \ker\bigl(H^{12}Z\longrightarrow H^{10}F\bigr)
       =\langle v_6^2\rangle.
 \tag{D.2.27e}
\]
These arguments are just the low path-spectral-sequence edge
calculations; they do not require a collapse in higher degrees.
The exact full-fiber relation
$\operatorname{Sq}^4e_6^*+\operatorname{Sq}^2J_1^*=0$
of \textup{D.2.24}, and suspension compatibility of squares,
now give \textup{(D.2.27b)}.

We next identify its class with the actual relative cup
functional, keeping one normalization for all maps.
Write $q:Z\to Y$. Choose actual cocycles $a_4,b_5$ and a
cochain $c_8$ on $Y$ with $dc_8=a_4\smile b_5$.
Since $q^*a_4$ is exact, choose $P_3$ with
$dP_3=q^*a_4$ and form
\[
 \Theta_8=[q^*c_8-P_3\smile q^*b_5]\in H^8Z.
 \tag{D.2.27f}
\]
Changing $P_3$ does not change this class: its difference
is a closed degree-three class multiplied by the exact
$q^*b_5$. Every $h:A\to Y$ lifts to $Z$, since
$h^*\alpha=0$ by \textup{D.2.12}. The actual cochain formula
of \textup{D.2.26} gives
\[
 \theta_h=\widetilde h^{,*}\Theta_8
 \quad\text{for every fiber-model lift }\widetilde h:A\to Z.
 \tag{D.2.27g}
\]
The independence just checked also handles the choice of lift.

Evaluate \textup{(D.2.27f)} on the counit
$\varepsilon:\Sigma^2F\to Z$. Its composite into $Y$ has
zero homology in degrees four and five by connectivity,
and in degree six because the bottom class $A_4=a^2$
has zero homology suspension. Thus \textup{(D.2.26j)} applies.
Its once-adjunct $\Sigma F\to\Omega Y$ has images on the
three degree-six fiber classes given by
\[
 d_6\longmapsto0,\qquad e_6\longmapsto ab+ba,
       \qquad f_6\longmapsto \text{the linear }f\text{ label}.
 \tag{D.2.27h}
\]
Here $a,b$ are the degree-three and degree-four letters in
$H_*\Omega Y$. The exact per-vector comparison
\textup{(D.2.26j)} therefore gives
$\sigma^2\Theta_8=e_6^*+t f_6^*$ for some $t\in\FF_2$.
The image of $\sigma^2q^*:H^8Y\to H^6F$ is precisely
$\langle f_6^*\rangle$. Changing $c_8$ by that one fixed
old cohomology class removes $t$. Since
$\sigma^2:H^8Z\to H^6F$ is an isomorphism, the resulting
universal class is exactly $\vartheta_8$.
This is a single change on $Y$, shared by every attaching map.

The homology of $A=Q\Sigma B$ is polynomial on primitive
generators: the unit labels are primitive because $\Sigma B$
is a suspension, and their Dyer--Lashof generators remain
primitive. Its dual Hopf algebra consequently has zero square
on every positive-dimensional cohomology class. This also
follows by pairing a square with the primitive-polynomial
coproduct: every nontrivial equal-factor term has even
coefficient. We use the opened free infinite-loop-space
description and Nishida formulas
\cite[Chapters I and III]{I_CohenLadaMay1976}.
Pulling \textup{(D.2.27b)} back along $\widetilde h$ gives
\[
 \operatorname{Sq}^4\widetilde h^{,*}\vartheta_8+
 \operatorname{Sq}^2\widetilde h^{,*}\lambda_{10}=0
                         \quad\text{in }H^{12}A.
 \tag{D.2.27i}
\]

We finish with five explicit witnesses, so no rank assertion
rests only on a computer calculation. Put
$A_0=\sigma a_2$, $B_0=\sigma b_3$, $C_5=\sigma c_4$,
$D_6=\sigma d_5$. Let $I_9,J_9$ be the unit classes dual
to $u^4,uv^2$, and let $G_8,H_8$ be the unit classes dual
to $u^2v,uw$. The complete basis of $H_8A$ is
\[
 Q^5A_0,\quad A_0C_5,\quad H_8,\quad G_8,\quad B_0^2.
 \tag{D.2.27j}
\]
For a base monomial $m$, write $\sigma(m^\vee)$ for its
unit homology dual, and put $t=\operatorname{Sq}^4w$.
The following degree-twelve classes all have zero second
square; their fourth squares are the displayed basis:
\[
\begin{array}{c|c|c}
 z\in H_{12}A&\operatorname{Sq}_*^2z&\operatorname{Sq}_*^4z\\ \hline
 Q^7C_5&0&Q^5A_0\\
 A_0(I_9+J_9)&0&A_0C_5\\
 \sigma((ut)^\vee)&0&H_8\\
 \sigma((u^4v)^\vee)+\sigma((u^3w)^\vee)&0&G_8\\
 D_6^2&0&B_0^2
\end{array}
 \tag{D.2.27k}
\]
Here the first row is the ordinary source Nishida identity:
the only nonzero term of $\operatorname{Sq}_*^4Q^7C_5$
is $Q^5\operatorname{Sq}_*^2C_5=Q^5A_0$, while the
leading coefficient for its second square is
$\binom52=0$ modulo two and $\beta C_5=0$.
The second row uses
$\operatorname{Sq}_*^2I_9=\operatorname{Sq}_*^2J_9=F_7$,
$\operatorname{Sq}_*^4I_9=C_5$, and
$\operatorname{Sq}_*^4J_9=0$.
The last row uses $\beta D_6=0$ and
$\operatorname{Sq}_*^2D_6=B_0$.
For both remaining rows, the complete relevant Cartan lists
in Serre's polynomial algebra are
\[
\begin{aligned}
 \operatorname{Sq}^2(u^3v)&=u^4v+uv^3+u^3w,&
 \operatorname{Sq}^2(u^2w)&=v^2w,\\
 \operatorname{Sq}^2(v^3)&=v^2w,&
 \operatorname{Sq}^2t&=0,\\
 \operatorname{Sq}^4(u^2v)&=u^4v+v^2w,&
 \operatorname{Sq}^4(uw)&=ut.
\end{aligned}
 \tag{D.2.27l}
\]
The degree-nine basis is $u^3v,u^2w,v^3,t$ and the
degree-seven basis is $u^2v,uw$, so these lists have no
omitted unit coefficient; compare
\cite[\S2, nos.~8--9]{I_Serre1953}.
All Dyer--Lashof calculations here occur inside the source;
no ordinary map is assumed to preserve those operations.

Pair \textup{(D.2.27i)} with each row of
\textup{(D.2.27k)}. It forces
$\widetilde h^{,*}\vartheta_8$ to vanish on the entire
basis \textup{(D.2.27j)}. Equations
\textup{(D.2.27g)--(D.2.27h)} prove
\textup{(D.2.27a)} with the one fixed normalization.
The tuple for any small simultaneous unsuspended family
is therefore zero. The exact common-cone criterion of
\textup{D.2.26} proves the asserted cup-product vanishing
and old-commutator survival.

For the mixed-family assertion, let $h:\Sigma^kA\to Y$
with $k\ge2$, lift it to $Z$, and let
$G:\Sigma^{k-2}A\to F$ be the double adjunct of that lift.
With exactly the universal normalization already fixed,
naturality of cohomology suspension gives
\[
 \theta_h=\widetilde h^{,*}\vartheta_8
       =\sigma^2(G^*e_6^*).
 \tag{D.2.27m}
\]
The entire functional $G^*e_6^*$ is zero for $k=2$ by
\textup{D.2.24}, and for $k=3$ and every $k\ge5$ by
\textup{D.2.25}. Thus every component in
\textup{(D.2.27a')} has zero Hopf functional with this one
choice of $c_8$. The simultaneous criterion applies to the
zero tuple, including all cross-attachment cancellations.
Nothing in this argument
identifies the new target after another cone with the
original $Y$, so it supplies no iteration claim.
\end{proof}
